% Options for packages loaded elsewhere
\PassOptionsToPackage{unicode}{hyperref}
\PassOptionsToPackage{hyphens}{url}
\PassOptionsToPackage{dvipsnames,svgnames,x11names}{xcolor}
%
\documentclass[
  11pt,
  letterpaper,
]{scrreprt}

\usepackage{amsmath,amssymb}
\usepackage{iftex}
\ifPDFTeX
  \usepackage[T1]{fontenc}
  \usepackage[utf8]{inputenc}
  \usepackage{textcomp} % provide euro and other symbols
\else % if luatex or xetex
  \usepackage{unicode-math}
  \defaultfontfeatures{Scale=MatchLowercase}
  \defaultfontfeatures[\rmfamily]{Ligatures=TeX,Scale=1}
\fi
\usepackage{lmodern}
\ifPDFTeX\else  
    % xetex/luatex font selection
    \setmainfont[]{TeX Gyre Pagella}
\fi
% Use upquote if available, for straight quotes in verbatim environments
\IfFileExists{upquote.sty}{\usepackage{upquote}}{}
\IfFileExists{microtype.sty}{% use microtype if available
  \usepackage[]{microtype}
  \UseMicrotypeSet[protrusion]{basicmath} % disable protrusion for tt fonts
}{}
\makeatletter
\@ifundefined{KOMAClassName}{% if non-KOMA class
  \IfFileExists{parskip.sty}{%
    \usepackage{parskip}
  }{% else
    \setlength{\parindent}{0pt}
    \setlength{\parskip}{6pt plus 2pt minus 1pt}}
}{% if KOMA class
  \KOMAoptions{parskip=half}}
\makeatother
\usepackage{xcolor}
\setlength{\emergencystretch}{3em} % prevent overfull lines
\setcounter{secnumdepth}{5}
% Make \paragraph and \subparagraph free-standing
\makeatletter
\ifx\paragraph\undefined\else
  \let\oldparagraph\paragraph
  \renewcommand{\paragraph}{
    \@ifstar
      \xxxParagraphStar
      \xxxParagraphNoStar
  }
  \newcommand{\xxxParagraphStar}[1]{\oldparagraph*{#1}\mbox{}}
  \newcommand{\xxxParagraphNoStar}[1]{\oldparagraph{#1}\mbox{}}
\fi
\ifx\subparagraph\undefined\else
  \let\oldsubparagraph\subparagraph
  \renewcommand{\subparagraph}{
    \@ifstar
      \xxxSubParagraphStar
      \xxxSubParagraphNoStar
  }
  \newcommand{\xxxSubParagraphStar}[1]{\oldsubparagraph*{#1}\mbox{}}
  \newcommand{\xxxSubParagraphNoStar}[1]{\oldsubparagraph{#1}\mbox{}}
\fi
\makeatother


\providecommand{\tightlist}{%
  \setlength{\itemsep}{0pt}\setlength{\parskip}{0pt}}\usepackage{longtable,booktabs,array}
\usepackage{calc} % for calculating minipage widths
% Correct order of tables after \paragraph or \subparagraph
\usepackage{etoolbox}
\makeatletter
\patchcmd\longtable{\par}{\if@noskipsec\mbox{}\fi\par}{}{}
\makeatother
% Allow footnotes in longtable head/foot
\IfFileExists{footnotehyper.sty}{\usepackage{footnotehyper}}{\usepackage{footnote}}
\makesavenoteenv{longtable}
\usepackage{graphicx}
\makeatletter
\def\maxwidth{\ifdim\Gin@nat@width>\linewidth\linewidth\else\Gin@nat@width\fi}
\def\maxheight{\ifdim\Gin@nat@height>\textheight\textheight\else\Gin@nat@height\fi}
\makeatother
% Scale images if necessary, so that they will not overflow the page
% margins by default, and it is still possible to overwrite the defaults
% using explicit options in \includegraphics[width, height, ...]{}
\setkeys{Gin}{width=\maxwidth,height=\maxheight,keepaspectratio}
% Set default figure placement to htbp
\makeatletter
\def\fps@figure{htbp}
\makeatother
% definitions for citeproc citations
\NewDocumentCommand\citeproctext{}{}
\NewDocumentCommand\citeproc{mm}{%
  \begingroup\def\citeproctext{#2}\cite{#1}\endgroup}
\makeatletter
 % allow citations to break across lines
 \let\@cite@ofmt\@firstofone
 % avoid brackets around text for \cite:
 \def\@biblabel#1{}
 \def\@cite#1#2{{#1\if@tempswa , #2\fi}}
\makeatother
\newlength{\cslhangindent}
\setlength{\cslhangindent}{1.5em}
\newlength{\csllabelwidth}
\setlength{\csllabelwidth}{3em}
\newenvironment{CSLReferences}[2] % #1 hanging-indent, #2 entry-spacing
 {\begin{list}{}{%
  \setlength{\itemindent}{0pt}
  \setlength{\leftmargin}{0pt}
  \setlength{\parsep}{0pt}
  % turn on hanging indent if param 1 is 1
  \ifodd #1
   \setlength{\leftmargin}{\cslhangindent}
   \setlength{\itemindent}{-1\cslhangindent}
  \fi
  % set entry spacing
  \setlength{\itemsep}{#2\baselineskip}}}
 {\end{list}}
\usepackage{calc}
\newcommand{\CSLBlock}[1]{\hfill\break\parbox[t]{\linewidth}{\strut\ignorespaces#1\strut}}
\newcommand{\CSLLeftMargin}[1]{\parbox[t]{\csllabelwidth}{\strut#1\strut}}
\newcommand{\CSLRightInline}[1]{\parbox[t]{\linewidth - \csllabelwidth}{\strut#1\strut}}
\newcommand{\CSLIndent}[1]{\hspace{\cslhangindent}#1}


% Wird für die Tabelle im Titelblatt der Experten verwendet:
% Array
\usepackage{array}
% Neue Definition für Tabelleneinträge
% linksbündig mit Breitenangabe
\newcolumntype{L}[1]{>{\raggedright\arraybackslash}p{#1}} 
% zentriert mit Breitenangabe
\newcolumntype{C}[1]{>{\centering\arraybackslash}p{#1}} 
% rechtsbündig mit Breitenangabe
\newcolumntype{R}[1]{>{\raggedleft\arraybackslash}p{#1}} 

\usepackage[a4paper, margin=3cm]{geometry}

% Titles take mainfont
\addtokomafont{disposition}{\rmfamily}





\makeatletter
\@ifpackageloaded{bookmark}{}{\usepackage{bookmark}}
\makeatother
\makeatletter
\@ifpackageloaded{caption}{}{\usepackage{caption}}
\AtBeginDocument{%
\ifdefined\contentsname
  \renewcommand*\contentsname{Inhaltsverzeichnis}
\else
  \newcommand\contentsname{Inhaltsverzeichnis}
\fi
\ifdefined\listfigurename
  \renewcommand*\listfigurename{Abbildungsverzeichnis}
\else
  \newcommand\listfigurename{Abbildungsverzeichnis}
\fi
\ifdefined\listtablename
  \renewcommand*\listtablename{Tabellenverzeichnis}
\else
  \newcommand\listtablename{Tabellenverzeichnis}
\fi
\ifdefined\figurename
  \renewcommand*\figurename{Abbildung}
\else
  \newcommand\figurename{Abbildung}
\fi
\ifdefined\tablename
  \renewcommand*\tablename{Tabelle}
\else
  \newcommand\tablename{Tabelle}
\fi
}
\@ifpackageloaded{float}{}{\usepackage{float}}
\floatstyle{ruled}
\@ifundefined{c@chapter}{\newfloat{codelisting}{h}{lop}}{\newfloat{codelisting}{h}{lop}[chapter]}
\floatname{codelisting}{Listing}
\newcommand*\listoflistings{\listof{codelisting}{Listingverzeichnis}}
\makeatother
\makeatletter
\makeatother
\makeatletter
\@ifpackageloaded{caption}{}{\usepackage{caption}}
\@ifpackageloaded{subcaption}{}{\usepackage{subcaption}}
\makeatother

\ifLuaTeX
\usepackage[bidi=basic]{babel}
\else
\usepackage[bidi=default]{babel}
\fi
\babelprovide[main,import]{ngerman}
\ifPDFTeX
\else
\babelfont{rm}[]{TeX Gyre Pagella}
\fi
% get rid of language-specific shorthands (see #6817):
\let\LanguageShortHands\languageshorthands
\def\languageshorthands#1{}
\ifLuaTeX
  \usepackage{selnolig}  % disable illegal ligatures
\fi
\usepackage{bookmark}

\IfFileExists{xurl.sty}{\usepackage{xurl}}{} % add URL line breaks if available
\urlstyle{same} % disable monospaced font for URLs
\hypersetup{
  pdftitle={Tragverhalten von Stahlbetontragwerken},
  pdfauthor={Pascal Gitz},
  pdflang={de},
  colorlinks=true,
  linkcolor={Black},
  filecolor={Maroon},
  citecolor={Blue},
  urlcolor={Blue},
  pdfcreator={LaTeX via pandoc}}


% TITELBLATT, VERSIONSTABELLE UND SELBSTSTÄNDIGKEITSERKLÄRUNG
%--------------------------------------------------------------------------------------------------------------------


\titlehead{\includegraphics[height=0.5cm]{../imgs/logos/logo-mse}\hfill\includegraphics[height=0.5cm]{../imgs/logos/logo-hslu-en-col} \\ }
\subject{MASTER OF SCIENCE IN ENGINEERING\\Vertiefungsmodul II}
\title{Tragverhalten von Stahlbetontragwerken}

\subtitle
{Modellvorstellungen zur nichtlinearen Verformungsberechnung}

%\thanks{
%Version 1.0
%\hfill \today
%\hfill Hun}


\date{\large Horw, Freitag, 17. Januar 2025}
\author{Pascal Gitz}

\publishers{
	\begin{table}[H]
		\centering
		\begin{tabular}{L{2cm} L{6cm}}
			Advisor: & Prof. FH, Dr. Daniel Heinzmann \\
			Experte: & Dr. Thomas Jäger \\
		\end{tabular}
	\end{table}
}
\begin{document}
\maketitle


Hiermit erkläre ich, dass ich die vorliegende Arbeit selbstständig angefertigt und keine anderen als die angegebenen Hilfsmittel verwendet habe. Sämtliche verwendeten Textausschnitte, Zitate oder Inhalte anderer Verfasser wurden ausdrücklich als solche gekennzeichnet.\\%
%
\\%
%
Horw, 14. Juni 2024 \hfill Pascal Gitz%

\vfill
%\begin{tabular}[h]{llcr} 
    %*Version 2.0 & - Definitives Exemplar & \today & MK \\ 
    %*Version 1.0 & - Prüfungsexemplar & 22. Januar 2021 & MK \\ 
%\end{tabular}\\

% Version 2.0 - Definitives Exemplar \hfill \today \quad \quad \quad \quad \quad MK\\
Version 1.0 - Prüfungsexemplar \hfill 14. Juni 2024 \quad \quad \quad \quad \quad PG\\
Version 0.9 - Entwurf \hfill 07. Juni 2024 \quad \quad \quad \quad \quad PG\\

\newpage

\chapter*{Kurzfassung}

Diese Arbeit ist Teil einer Arbeitsserie, bestehend aus dem Vertiefungsmodul I, dem Vertiefungsmodul II und der Masterthesis. Der Kern dieser Serie ist die Beschreibung von Verformungen im Stahlbetonbau basierend auf mechanischen Modellen. In der ersten Vertiefungsarbeit wurden die Hintergründe der Modelle beleuchtet und mittels Handrechnungen an Versuchen verifiziert. In dieser Arbeit, dem Vertiefungsmodul II, werden die Ansätze aus der Vorarbeit mit kommerziellen Finite-Element-Tools (FE-Tools) kombiniert. Es werden die Softwares AxisVM-X7, RFEM-6 und IDEAStatiCa verwendet.

Eine Modellvorstellung, das Federmodell, wird aufgezeigt. Dabei werden biege- und schubstarre Stäbe in regelmässigen Abständen durch Federn gekoppelt, in denen beliebig nichtlineare Steifigkeitsbeziehungen hinterlegt werden können. Dies ermöglicht die Bestimmung nichtlinearer Verformungen. Die Steifigkeiten werden anhand der Momenten-Krümmungs-Beziehung und anhand der Steifigkeit der Schubbewehrung ermittelt. Das Modell wird an einem Einführungsbeispiel illustriert, wobei analytische Lösungen mit den FE-Ergebnissen verglichen werden. Diese zeigen, dass die Ergebnisse nahezu deckungsgleich und zufriedenstellend sind, mit minimalen numerischen Abweichungen, die auf numerische Approximationen zurückzuführen sind.

Im Anschluss wird das Modell verifiziert. In diesem Kapitel werden die Versuche aus der Vorarbeit mit dem Federmodell nachgerechnet, um identische Ergebnisse wie bei den Handrechnungen zu erzielen. An einem Dreipunktbiegeversuch werden Schub- und Biegeverformungen basierend auf nichtlinearen Stoffgesetzen ermittelt und mit den gemessenen Versuchsresultaten verglichen. Die Ergebnisse sind zufriedenstellend. An einem Vierpunktbiegeversuch werden ebenfalls Biege- und Schubverformungen bestimmt und die Anwendung des Versatzmasses aufgezeigt. Auch hier werden die Modellresultate mit den Versuchsresultaten verglichen, wobei Abweichungen erkennbar sind, die knapp akzeptabel sind. Diese Unterschiede werden auf Unschärfen in den angewendeten Stoffgesetzen zurückgeführt. Im Vergleich mit den Handrechnungen aus der Vorarbeit zeigt das Federmodell jedoch identische Ergebnisse. Die Modellverifizierung wird abschliessend als erfolgreich betrachtet.

Im Kapitel der Modellanwendung wird das Federmodell an einem vorgespannten Träger angewendet. Das Kapitel beginnt mit einer detaillierten Versuchsbeschreibung, gefolgt von der Modellbildung. Die Interpretation der Vorspannung als Einwirkung sowie der Lastfall der Pressenkräfte und das Eigengewicht werden beschrieben. Der Hauptaspekt der Analyse ist die Ermittlung der Momenten-Krümmungs-Beziehung anhand nichtlinearer Stoffgesetze. Mittels einer Querschnittsanalysesoftware werden entlang der Stabachse 136 unterschiedliche Querschnitte analysiert. Die numerischen Ergebnisse werden erfolgreich durch Handrechnungen verifiziert. Die Resultate des Modells werden mit den Versuchsmessungen verglichen, wobei ein annähernd deckungsgleicher Verlauf festgestellt wird. Lediglich im Bereich der Traglast sind Unterschiede erkennbar. Das Kapitel endet mit einem Vergleich des Federmodells mit der Compatible Stress Field Method (CSFM). Die CSFM zeigt ähnliche Resultate wie das Federmodell.

\renewcommand*\contentsname{Inhaltsverzeichnis}
{
\hypersetup{linkcolor=}
\setcounter{tocdepth}{1}
\tableofcontents
}
\listoffigures

\bookmarksetup{startatroot}

\chapter{Einleitung}\label{einleitung}

Bereits 1935 wurden in der damaligen SIA 112 Anforderungen an die
Durchbiegung von Stahlbetonbauteilen gestellt. Seither wurde die
Gebrauchstauglichkeit, insbesondere die Verformungen von Tragwerken, zu
einer zentralen Bemessungsgrösse. Trotz dieser elementaren Rolle im
Bemessungsablauf werden Verformungen im Stahlbetonbau in der heutigen
Baupraxis häufig nur mit groben Abschätzungen bewertet. Solche Ansätze
können zu ungenauen Prognosen führen, die sowohl die Sicherheit als auch
die Wirtschaftlichkeit von Bauprojekten beeinträchtigen. Detaillierte
Verformungsberechnungen sind in vielen modernen Statiksoftwares
implementiert und frei anwendbar. Die theoretischen Hintergründe sind
jedoch selten geläufig, und die Tools werden meist als
\emph{``Blackbox''} behandelt.

Das Ziel dieser Arbeitsserie, bestehend aus dem Vertiefungsmodul I
{[}\citeproc{ref-gitz_ansatze_2024}{1}{]}, Vertiefungsmodul II und der
Masterthesis, ist es daher, dem Leser mechanische Modelle zur Bestimmung
von Verformungen umfassend aufzuzeigen. In der Vorarbeit
{[}\citeproc{ref-gitz_ansatze_2024}{1}{]} wurden die mechanischen
Modelle detailliert beschrieben und anhand von experimentellen Versuchen
verifiziert, wobei mehrheitlich auf den Einsatz kommerzieller
Finite-Elemente-Tools (FE-Tools) verzichtet wurde. Der Fokus dieser
Arbeit liegt nun auf der Integration von FE-Tools mit den theoretischen
Grundlagen, um praktisch anwendbare Lösungen zu entwickeln. Verwendet
wird hierbei die 3D Statiksoftware AxisVM-X7 und RFEM-6, sowie die 2D
Detailstatiksoftware IDEAStatiCa. Sämtliche Analysen beziehen sich auf
Stabtragwerke.

Diese Arbeit lässt sich in drei Kernthemen unterteilen. Zuerst wird die
Modellvorstellung zur numerischen Integration kinematischer Grössen
(primär Krümmung und Schiebung) beschrieben und anhand eines
Einführungsbeispiels veranschaulicht, um die theoretischen Grundlagen
klar darzustellen. Anschliessend wird die Modellvorstellung an den
bereits gut verstandenen Versuchen aus dem Vertiefungsmodul I
verifiziert, um die Zuverlässigkeit und Genauigkeit des Modells zu
bestätigen. Abschliessend wird die Anwendung an einem vorgespannten
Träger demonstriert, um die praktische Relevanz und Anwendbarkeit des
entwickelten Modells zu zeigen. Die Arbeit wird mit einem Fazit
abgeschlossen, das eine differenzierte Betrachtung zwischen Rückblick
und Ausblick bietet und damit sowohl die erreichten Ergebnisse
zusammenfasst als auch Erweiterungen aufzeigt.

Durch diese strukturierte Herangehensweise soll ein tieferes Verständnis
für das komplexe Tragverhalten von Stahlbetontragwerken geschaffen
werden. Zudem soll eine zuverlässige Grundlage zur praktischen Anwendung
im konstruktiven Ingenieurbau gelegt werden.

\phantomsection\label{refs}
\begin{CSLReferences}{0}{1}
\bibitem[\citeproctext]{ref-gitz_ansatze_2024}
\CSLLeftMargin{1. }%
\CSLRightInline{Gitz P (2024) Ansätze zur {Verformungsberechnung}. HSLU
Technik \& Architektur}

\end{CSLReferences}




\end{document}
